
\documentclass{article}
\usepackage{colortbl}
\usepackage{makecell}
\usepackage{multirow}
\usepackage{supertabular}

\begin{document}

\newcounter{utterance}

\centering \large Interaction Transcript for game `eqbench', experiment `eq\_bench', episode 59 with reflexion{-}Qwen3.5{-}9B.
\vspace{24pt}

{ \footnotesize  \setcounter{utterance}{1}
\setlength{\tabcolsep}{0pt}
\begin{supertabular}{c@{$\;$}|p{.15\linewidth}@{}p{.15\linewidth}p{.15\linewidth}p{.15\linewidth}p{.15\linewidth}p{.15\linewidth}}
    \# & \multicolumn{2}{c}{Player} && \multicolumn{2}{c}{Game Master} \\
    \hline

    \theutterance \stepcounter{utterance}  
    & & & \multicolumn{4}{p{0.6\linewidth}}{
        \cellcolor[rgb]{0.9,0.9,0.9}{
            \makecell[{{p{\linewidth}}}]{
                \texttt{\tiny{[P1$\langle$GM]}}
                \texttt{You have attempted this same task before and did not succeed. These are the plans you wrote after those attempts.} \\
\texttt{Trial \#1: New plan: I failed because I assigned a high score to "Desperation" (8) and "Remorse" (7) without sufficiently analyzing the dialogue's tone, which is characterized by Tyler's aggression and Claudia's defensive, resigned acceptance rather than active pleading or deep guilt. "Desperation" is low because Claudia is not begging; she is stating facts ("I wanted to see you") despite his accusations. "Remorse" is also low because she explicitly denies being a burden and frames her past actions as necessary ("I had to"), showing a lack of regret. "Indifference" is incorrect as she is clearly emotionally engaged. The correct primary emotion is "Embarrassment," driven by her stammering ("I..."), her defensive denial ("That's not fair"), and her inability to counter his harsh words, indicating she feels ashamed of her past failure to protect him or explain herself better. Next time, I will prioritize "Embarrassment" as the highest score (9{-}10), assign a low score to "Desperation" (1{-}2) due to her lack of pleading, and a low score to "Remorse" (2{-}3) due to her rationalization of her past actions.} \\
\\ 
\texttt{Trial \#2: New plan: I failed because I assigned a high score to "Embarrassment" (8) and "Remorse" (3) while underestimating "Desperation" and "Indifference." Upon re{-}evaluating the dialogue, Tyler's accusations ("lies," "abandoned me," "pretend everything is okay") are severe and accusatory, causing Claudia to break down with fragmented speech ("I...I wanted to see you," "I..."). This stammering and her inability to defend herself against such harsh judgment indicate a state of high anxiety and fear of rejection, which aligns much more strongly with **Desperation** than simple embarrassment. While she denies being a burden, her tone is defensive and pleading rather than purely ashamed. "Indifference" is incorrect as she is clearly distressed, and "Remorse" is low because she rationalizes her past ("I had to") rather than expressing deep regret. Next time, I will assign the highest score (9{-}10) to **Desperation** due to her emotional collapse under pressure, a low score (1{-}2) to **Embarrassment** as it is secondary to her fear, a low score (2{-}3) to **Remorse** due to her rationalization, and a zero to **Indifference**.} \\
\\ 
\\ 
\\ 
\texttt{Your task is to predict the likely emotional responses of a character in this dialogue:} \\
\\ 
\texttt{Claudia: Tyler, you look well. Teaching suits you.} \\
\texttt{Tyler: Yeah, well, someone's gotta do the dirty work.} \\
\texttt{Claudia: I didn't mean it like that. I...} \\
\texttt{Tyler: I know what you meant. You always do, don't you?} \\
\texttt{Claudia: Not always, no. But I try to understand, at least.} \\
\texttt{Tyler: Understand? Like you understood when you left me alone?} \\
\texttt{Claudia: That's not fair, Ty. I had to...} \\
\texttt{Tyler: What? Chase your dreams? Leave your past behind?} \\
\texttt{Claudia: I never left you behind, Tyler. You were always...} \\
\texttt{Tyler: Always what? An afterthought? A burden?} \\
\texttt{Claudia: No, never a burden. You were my responsibility.} \\
\texttt{Tyler: Responsibility? Is that why you're here now? After all these years?} \\
\texttt{Claudia: I'm here because I...I wanted to see you.} \\
\texttt{Tyler: Or maybe you're here because you need something. Like always.} \\
\texttt{Claudia: That's not true, Tyler. I...} \\
\texttt{Tyler: Save it, Claudia. I've heard enough of your lies.} \\
\texttt{Claudia: I'm not lying, Tyler. I...} \\
\texttt{Tyler: You abandoned me, Claudia. You can't just waltz back into my life and pretend everything is okay.} \\
\texttt{{[}End dialogue{]}} \\
\\ 
\texttt{At the end of this dialogue, Claudia would feel...} \\
\texttt{Embarrassment} \\
\texttt{Desperation} \\
\texttt{Indifference} \\
\texttt{Remorse} \\
\\ 
\texttt{Give each of these possible emotions a score from 0{-}10 for the relative intensity that they are likely to be feeling each.} \\
\\ 
\texttt{You must output in the following format, including headings (of course, you should give your own scores), with no additional commentary:} \\
\\ 
\texttt{Embarrassment: <score>} \\
\texttt{Desperation: <score>} \\
\texttt{Indifference: <score>} \\
\texttt{Remorse: <score>} \\
\\ 
\\ 
\texttt{{[}End of answer{]}} \\
\\ 
\texttt{Remember: zero is a valid score, meaning they are likely not feeling that emotion. You must score at least one emotion > 0.} \\
\\ 
\texttt{Your answer:} \\
            }
        }
    }
     \\ \\

    \theutterance \stepcounter{utterance}  
    & \multicolumn{4}{p{0.6\linewidth}}{
        \cellcolor[rgb]{0.561,0.737,0.561}{
            \makecell[{{p{\linewidth}}}]{
                \texttt{\tiny{[P1$\rangle$GM]}}
                \texttt{Embarrassment: 2} \\
\texttt{Desperation: 9} \\
\texttt{Indifference: 0} \\
\texttt{Remorse: 3} \\
            }
        }
    }
    & & \\ \\

    \theutterance \stepcounter{utterance}  
    & & & \multicolumn{2}{p{0.3\linewidth}}{
        \cellcolor[rgb]{0.9,0.9,0.9}{
            \makecell[{{p{\linewidth}}}]{
                \texttt{\tiny{[GM$|$GM]}}
                \texttt{Embarrassment: 2} \\
\texttt{Desperation: 9} \\
\texttt{Indifference: 0} \\
\texttt{Remorse: 3} \\
            }
        }
    }
    & & \\ \\

    \theutterance \stepcounter{utterance}  
    & & & \multicolumn{2}{p{0.3\linewidth}}{
        \cellcolor[rgb]{0.9,0.9,0.9}{
            \makecell[{{p{\linewidth}}}]{
                \texttt{\tiny{[GM$|$GM]}}
                \texttt{Embarrassment: 6} \\
\texttt{Desperation: 7} \\
\texttt{Indifference: 0} \\
\texttt{Remorse: 8} \\
            }
        }
    }
    & & \\ \\

    \theutterance \stepcounter{utterance}  
    & & & \multicolumn{2}{p{0.3\linewidth}}{
        \cellcolor[rgb]{0.9,0.9,0.9}{
            \makecell[{{p{\linewidth}}}]{
                \texttt{\tiny{[GM$|$GM]}}
                \texttt{game\_result = LOSE} \\
            }
        }
    }
    & & \\ \\

    \theutterance \stepcounter{utterance}  
    & & \multicolumn{4}{p{0.6\linewidth}}{
        \cellcolor[rgb]{0.95,0.95,0.95}{
            \makecell[{{p{\linewidth}}}]{
                \texttt{\tiny{[P1$\leftrightarrow$P]}}
                \texttt{Post{-}episode Reflexion (not sent to GM):} \\
\texttt{New plan: I failed because my previous attempts oscillated between prioritizing "Desperation" and "Embarrassment" without fully accounting for the specific nuance of Claudia's defensive rationalization versus her emotional collapse. While her stammering ("I...") indicates distress, her repeated attempts to justify her actions ("I had to...", "You were my responsibility") suggest she is trying to manage the situation through explanation rather than pure panic. However, Tyler's final line ("Save it... I've heard enough of your lies") completely shuts down her ability to explain, leaving her with no recourse but to accept his judgment. This specific dynamic—where a character tries to defend a past action and is met with an absolute refusal to listen—creates a state of high **Desperation** (fear of total rejection) and significant **Remorse** (the realization that her explanations are futile and she is being seen as a liar). "Embarrassment" is present but secondary to the fear of abandonment, and "Indifference" is incorrect. The correct strategy is to assign the highest score (9{-}10) to **Desperation** due to the immediate threat of being cut off, a high score (8{-}9) to **Remorse** because her defenses have crumbled under his accusations, a moderate score (4{-}5) to **Embarrassment** as a reaction to his harsh words, and zero to **Indifference**.} \\
            }
        }
    }
    & \\ \\

\end{supertabular}
}

\end{document}
